\documentclass[12pt,a4paper,twoside]{ctexbook}

\usepackage{amsmath,amssymb}
\usepackage{graphicx}
\usepackage{booktabs}
\usepackage{longtable}
\usepackage{hyperref}

\title{基于机器学习的城市交通流预测方法研究}

\begin{document}

% ===== 中文摘要 =====
\begin{abstract}
随着城市化进程的加速，城市交通拥堵问题日益严重，精准的交通流预测对智能交通系统建设具有重要意义。本文系统梳理了交通流预测领域的研究现状，分析了传统统计模型与深度学习模型的优缺点。在此基础上，提出了一种融合注意力机制的时空图卷积网络模型。该模型通过图卷积网络捕捉路网拓扑结构中的空间依赖关系，利用门控循环单元提取时间序列特征，并通过多头注意力机制增强长距离依赖建模能力。实验结果表明，所提模型在PeMSD4和PeMSD8数据集上的预测精度均优于基线模型，平均绝对误差降低幅度达到百分之八以上。此外，本文还对模型的可解释性进行了分析，探讨了不同类型的注意力头在时空建模中的实际作用。全文包含六章内容，系统地阐述了从问题提出到模型构建、实验验证及结论总结的完整研究过程。

交通流预测是智能交通系统中的核心技术问题。随着传感器技术的普及和海量交通数据的积累，数据驱动的预测方法逐渐成为主流。然而，现有方法在建模复杂的时空依赖关系时仍面临诸多挑战。本文的创新点主要体现在三个方面：第一，设计了一种自适应邻接矩阵，能够动态调整路网节点之间的连接权重；第二，提出了时空联合嵌入表示，将时间和空间信息融合到一个统一的向量空间中；第三，引入了课程学习策略，通过由易到难的训练方式提升模型的泛化性能。
\end{abstract}

\keywords{交通流预测，时空图卷积网络，注意力机制，深度学习，智能交通系统}

% ===== 英文摘要 =====
\begin{abstract}[english]
With the acceleration of urbanization, urban traffic congestion has become increasingly severe. Accurate traffic flow prediction is of great significance for the construction of intelligent transportation systems. This thesis systematically reviews the current state of research in the field of traffic flow prediction and analyzes the strengths and weaknesses of traditional statistical models and deep learning models. On this basis, a spatio-temporal graph convolutional network model integrating attention mechanisms is proposed. The model captures spatial dependencies in road network topology through graph convolutional networks, extracts temporal features using gated recurrent units, and enhances long-range dependency modeling capability through multi-head attention mechanisms. Experimental results demonstrate that the proposed model outperforms baseline models on both the PeMSD4 and PeMSD8 datasets, with a reduction in mean absolute error exceeding eight percent. Furthermore, this thesis analyzes the interpretability of the model and explores the practical roles of different types of attention heads in spatio-temporal modeling. The thesis consists of six chapters, systematically presenting the complete research process from problem formulation to model construction, experimental validation, and conclusion summarization.

Traffic flow prediction is a core technical problem in intelligent transportation systems. With the widespread deployment of sensors and the accumulation of massive traffic data, data-driven prediction methods have gradually become mainstream. However, existing methods still face many challenges in modeling complex spatio-temporal dependencies. The innovations of this thesis are mainly reflected in three aspects: first, an adaptive adjacency matrix is designed that can dynamically adjust the connection weights between road network nodes; second, a spatio-temporal joint embedding representation is proposed, integrating temporal and spatial information into a unified vector space; third, a curriculum learning strategy is introduced to enhance the generalization performance of the model through an easy-to-hard training approach.
\end{abstract}

\keywords[english]{traffic flow prediction, spatio-temporal graph convolutional network, attention mechanism, deep learning, intelligent transportation system}

\tableofcontents

% ===== 第1章 引言 =====
\chapter{引言}
\section{研究背景与意义}
随着我国城镇化率的持续提升，城市人口规模不断扩大，机动车保有量快速增长。根据公安部交通管理局的统计数据，截至2025年底，全国机动车保有量已突破四亿辆。城市交通拥堵不仅造成巨大的经济损失，还带来了严重的环境污染问题。因此，发展智能交通系统成为解决城市交通问题的重要途径。

智能交通系统的核心功能之一是交通流预测，即基于历史和当前的交通状态信息，预测未来一段时间内的交通流量、速度和密度等参数。准确的交通流预测能够为交通管理部门提供决策支持，帮助出行者合理规划路线，从而有效缓解交通拥堵。然而，交通流具有复杂的时空依赖特性，对其进行精准建模仍然是一个具有挑战性的研究课题。

\section{国内外研究现状}
\subsection{传统交通流预测方法}
早期的交通流预测研究主要采用统计方法和传统机器学习方法。历史平均模型是其中最简单的一类方法，但无法捕捉交通流的动态变化特征。自回归积分滑动平均模型及其变体能够较好地建模时间序列数据的线性特征，但对非线性关系的拟合能力有限。支持向量回归和k近邻等传统机器学习方法在一定程度上弥补了线性模型的不足，但由于需要手工设计特征，模型的性能在很大程度上依赖于特征工程的质量。

\subsection{深度学习预测方法}
近年来，深度学习技术在交通流预测领域得到了广泛应用。循环神经网络及其变体如长短期记忆网络和门控循环单元能够有效捕捉时序数据的长期依赖关系，成为交通流预测的主流模型。卷积神经网络在提取空间特征方面的优势也被用于处理交通数据的空间相关性。

图卷积网络的出现为解决交通路网的非欧几里得空间结构提供了新的思路。通过将交通路网建模为图结构，图卷积网络能够自然地捕捉节点之间的连接关系。在此基础上，多种时空图网络模型被提出，如扩散卷积循环神经网络、时空图卷积网络和基于注意力机制的时空图网络等。

\section{本文的主要工作与创新点}
本文围绕城市交通流预测问题展开研究，主要工作和创新点包括以下方面：提出了一种融合多头注意力机制的时空图卷积网络模型，设计了自适应邻接矩阵以动态建模路网空间依赖关系，引入了时空联合嵌入表示方法，并采用课程学习策略提升模型训练效果。

\section{论文的结构安排}
本文共分为六章。第一章介绍研究背景和意义，综述国内外研究现状，明确本文的研究内容和创新点。第二章介绍相关的理论基础，包括图卷积网络、循环神经网络和注意力机制。第三章详细阐述本文提出的模型结构。第四章介绍实验设计和结果分析。第五章对全文进行总结，并展望未来的研究方向。此外，本文还包含参考文献、附录和致谢等部分。

% ===== 第2章 相关理论基础（含图表） =====
\chapter{相关理论基础}
\section{图卷积网络}
图卷积网络是对传统卷积神经网络在处理非欧几里得数据方面的延伸。在图结构中，每个节点的邻域结构各不相同，因此无法直接应用标准卷积操作。谱域图卷积和空域图卷积是两类主要的图卷积方法。谱域方法利用图的拉普拉斯矩阵进行傅里叶变换，在频域中定义卷积运算。空域方法则直接在节点的邻域上定义卷积操作，通过聚合邻居节点的信息来更新中心节点的表示。

\begin{figure}[htbp]
\centering
\rule{0.75\textwidth}{4cm}
\caption{图卷积网络的信息聚合过程示意图}
\label{fig:gcn}
\end{figure}

图卷积网络的信息聚合过程可以通过邻接矩阵和特征矩阵的乘积来实现。图~\ref{fig:gcn} 展示了图卷积网络中节点信息聚合的基本流程。输入层接收图中所有节点的特征向量，经过多层图卷积操作后，每个节点的表示都融合了其局部邻域的结构信息和属性信息。

\section{循环神经网络与门控机制}
长短期记忆网络通过引入遗忘门、输入门和输出门三个门控单元来调节信息的流动，从而有效缓解了传统循环神经网络中的梯度消失问题。门控循环单元则对长短期记忆网络进行了简化，将遗忘门和输入门合并为更新门，减少了参数量，在保持相当性能的同时提高了计算效率。

\begin{figure}[htbp]
\centering
\rule{0.65\textwidth}{3.5cm}
\caption{门控循环单元的内部结构}
\label{fig:gru}
\end{figure}

图~\ref{fig:gru} 展示了门控循环单元的内部结构。其中重置门控制前一时间步状态信息中被保留的比例，更新门决定当前时间步状态更新的幅度。这两个门控信号共同决定了序列信息在网络中的传输方式。

\section{注意力机制}
注意力机制最早在自然语言处理领域被提出，用于解决序列到序列模型中的长距离依赖问题。其核心思想是计算输入序列中各元素之间的关联权重，使模型能够自适应地关注输入中的重要部分。多头注意力机制通过并行执行多次注意力计算，使模型能够在不同的表示子空间中捕获不同的特征模式。

\begin{figure}[htbp]
\centering
\rule{0.7\textwidth}{3cm}
\caption{多头注意力机制的计算流程}
\label{fig:attention}
\end{figure}

缩放点积注意力是多头注意力机制的核心组件，其计算过程如图~\ref{fig:attention} 所示。通过引入缩放因子，避免了点积值过大导致的梯度不稳定问题。

\section{交通流数据特征分析}
\begin{table}[htbp]
\centering
\caption{PeMS数据集统计信息}
\label{tab:dataset}
\begin{tabular}{lcccc}
\toprule
数据集 & 传感器数量 & 时间步数 & 时间间隔 & 缺失率 \\
\midrule
PeMSD4 & 307 & 16992 & 5分钟 & 3.18\% \\
PeMSD7 & 883 & 28224 & 5分钟 & 2.46\% \\
PeMSD8 & 170 & 17856 & 5分钟 & 2.73\% \\
\bottomrule
\end{tabular}
\end{table}

表~\ref{tab:dataset} 总结了三个常用的交通流预测基准数据集的基本统计信息。这些数据集均来自加州交通局的性能测量系统，覆盖了不同规模和复杂度的交通路网。

\begin{table}[htbp]
\centering
\caption{不同预测步长下的模型性能对比（MAE）}
\label{tab:performance}
\begin{tabular}{lcccc}
\toprule
模型 & 15分钟 & 30分钟 & 45分钟 & 60分钟 \\
\midrule
HA & 3.82 & 3.82 & 3.82 & 3.82 \\
LSTM & 2.91 & 3.24 & 3.57 & 3.89 \\
GRU & 2.87 & 3.18 & 3.49 & 3.81 \\
STGCN & 2.64 & 2.98 & 3.30 & 3.62 \\
本文模型 & 2.41 & 2.73 & 3.02 & 3.31 \\
\bottomrule
\end{tabular}
\end{table}

从表~\ref{tab:performance} 可以看出，本文提出的模型在所有预测步长下均取得了最优的平均绝对误差。随着预测步长的增加，各模型的预测误差均有所上升，但本文模型的性能衰减最为平缓，表明其对长时预测具有较好的鲁棒性。

% ===== 第3章 模型设计与实现（含公式） =====
\chapter{模型设计与实现}
\section{问题定义}
交通流预测问题可以形式化为一个时空序列预测任务。给定一个交通路网中 $N$ 个传感器在 $T$ 个历史时间步内观测到的交通数据，目标是预测未来 $H$ 个时间步的交通状态。这一过程可以用数学表达式进行描述。

\section{空间依赖建模}
为了捕捉交通路网中的空间依赖关系，本文采用图卷积网络对路网拓扑结构进行建模。第一层图卷积操作可以表示为：
\begin{equation}
H^{(1)} = \sigma\left(\tilde{D}^{-\frac{1}{2}} \tilde{A} \tilde{D}^{-\frac{1}{2}} X W^{(0)}\right)
\label{eq:gcn1}
\end{equation}
其中 $\tilde{A} = A + I_N$ 表示添加了自环的邻接矩阵，$\tilde{D}$ 为对应的度矩阵，$X$ 为输入特征矩阵，$W^{(0)}$ 为可学习的权重参数，$\sigma$ 为非线性激活函数。

本文进一步引入了自适应邻接矩阵，其计算方式为：
\begin{equation}
A_{adp} = \text{SoftMax}\left(\text{ReLU}(E_1 E_2^T)\right)
\label{eq:adp}
\end{equation}
其中 $E_1, E_2 \in \mathbb{R}^{N \times d}$ 是两个可学习的节点嵌入矩阵，$d$ 为嵌入维度。式~\eqref{eq:adp} 通过嵌入向量的内积计算节点间的关联强度，并使用SoftMax函数进行归一化。

\section{时间依赖建模}
对于时间维度的特征提取，本文采用门控循环单元对交通流序列进行建模：
\begin{equation}
z_t = \sigma\left(W_z \cdot [h_{t-1}, x_t] + b_z\right)
\label{eq:gru_z}
\end{equation}
\begin{equation}
r_t = \sigma\left(W_r \cdot [h_{t-1}, x_t] + b_r\right)
\label{eq:gru_r}
\end{equation}
其中 $z_t$ 和 $r_t$ 分别表示更新门和重置门的输出，$h_{t-1}$ 为上一时间步的隐藏状态。

最终的隐藏状态更新公式为：
\begin{equation}
h_t = (1 - z_t) \odot h_{t-1} + z_t \odot \tanh\left(W_h \cdot [r_t \odot h_{t-1}, x_t] + b_h\right)
\label{eq:gru_h}
\end{equation}
通过式~\eqref{eq:gru_h} 的递推计算，门控循环单元能够有效地在时间维度上传播信息，捕获交通流序列中的长期依赖关系。

\section{损失函数设计}
模型的训练目标是最小化预测值与真实值之间的差距。本文采用Huber损失函数：
\begin{equation}
\mathcal{L} = \frac{1}{N \cdot H} \sum_{i=1}^{N} \sum_{j=1}^{H} \ell\left(\hat{y}_{i,j}, y_{i,j}\right)
\label{eq:loss}
\end{equation}
其中 $\ell(\cdot)$ 为Huber损失函数，该损失函数在误差较小时等价于均方误差，在误差较大时等价于绝对误差，因此对异常值具有更好的鲁棒性。

% ===== 参考文献 =====
\begin{thebibliography}{99}

\bibitem{arellano1991}
Arellano M, Bond S. Some Tests of Specification for Panel Data: Monte Carlo Evidence and an Application to Employment Equations[J]. Review of Economic Studies, 1991, 58(2): 277-297.

\bibitem{vaswani2017}
Vaswani A, Shazeer N, Parmar N, et al. Attention Is All You Need[C]. Advances in Neural Information Processing Systems, 2017: 5998-6008.

\bibitem{kipf2017}
Kipf T N, Welling M. Semi-Supervised Classification with Graph Convolutional Networks[C]. International Conference on Learning Representations, 2017.

\bibitem{hochreiter1997}
Hochreiter S, Schmidhuber J. Long Short-Term Memory[J]. Neural Computation, 1997, 9(8): 1735-1780.

\bibitem{li2018}
Li Y, Yu R, Shahabi C, et al. Diffusion Convolutional Recurrent Neural Network: Data-Driven Traffic Forecasting[C]. International Conference on Learning Representations, 2018.

\bibitem{yu2018}
Yu B, Yin H, Zhu Z. Spatio-Temporal Graph Convolutional Networks: A Deep Learning Framework for Traffic Forecasting[C]. International Joint Conference on Artificial Intelligence, 2018: 3634-3640.

\bibitem{guo2019}
Guo S, Lin Y, Feng N, et al. Attention Based Spatial-Temporal Graph Convolutional Networks for Traffic Flow Forecasting[C]. AAAI Conference on Artificial Intelligence, 2019, 33: 922-929.

\bibitem{王明2020}
王明, 张伟, 李强. 基于深度学习的城市交通流预测综述[J]. 计算机学报, 2020, 43(3): 547-569.

\bibitem{陈静2021}
陈静, 刘洋. 时空图卷积网络在交通预测中的应用研究[J]. 自动化学报, 2021, 47(8): 1817-1830.

\bibitem{zheng2020}
Zheng C, Fan X, Wang C, et al. GMAN: A Graph Multi-Attention Network for Traffic Prediction[C]. AAAI Conference on Artificial Intelligence, 2020, 34(1): 1234-1241.

\end{thebibliography}

% ===== 附录 =====
\appendix

\chapter{实验环境配置}
本文所有实验均在如下软硬件环境下进行。操作系统为Ubuntu 22.04 LTS，处理器为Intel Xeon Gold 6248R，主频为3.0GHz，内存为256GB DDR4。GPU为NVIDIA Tesla V100，显存容量为32GB。深度学习框架采用PyTorch 2.0.1，CUDA版本为11.8。

模型训练过程中使用Adam优化器，初始学习率设定为零点零零一。批量大小根据数据集规模分别设定为三十二或六十四。所有模型均训练一百个周期，并采用早停策略，当验证集损失连续十个周期未下降时终止训练。

\chapter{补充实验结果}
本节给出模型在更多评价指标下的详细性能对比。除了主实验中报告的平均绝对误差之外，还包括均方根误差和平均绝对百分比误差。从补充结果来看，本文模型在所有数据集和所有评价指标上均表现出较好的一致性和稳定性。

为了检验模型的统计显著性，本文还进行了Wilcoxon符号秩检验。检验结果表明，本文模型与次优基线模型之间的性能差异在百分之九十五的置信水平下具有统计显著性，验证了所提方法的有效性。

% ===== 后记 =====
\acknowledgments{
时光荏苒，博士求学生涯即将画上句号。回首过去的几年，感慨万千。

首先，衷心感谢我的导师张教授。在博士学习期间，张老师不仅在学术研究上给予我悉心指导，还在生活中给予我无私关怀。张老师严谨的治学态度、开阔的学术视野和谦逊的人格魅力，将永远是我学习的榜样。

感谢实验室的各位同门，与你们一起度过的时光是我人生中最宝贵的财富。特别感谢李明同学在实验设计和论文写作过程中提供的帮助和建议。感谢我的室友王磊，在无数个挑灯夜读的夜晚，我们互相鼓励、共同进步。

感谢我的父母，你们多年的养育之恩和默默支持，是我不断前行的动力。感谢我的女友陈婉莹，在我科研遇到瓶颈时的耐心陪伴和鼓励，让我能够以积极的心态面对挑战。

最后，感谢百忙之中参与本文评审和答辩的各位专家老师，你们的宝贵意见将使本文更加完善。

博士生涯的结束不是终点，而是新征程的起点。我将带着在这里收获的知识和友谊，继续探索未知的世界。
}

\end{document}
